\section{Three inference servers, and why llama.cpp}\label{sec:servers}

I tried three servers before settling on one: llama.cpp with GGUF builds, vLLM with NVIDIA's NVFP4
checkpoint, and SGLang. This section reports what each could do on this hardware. All of it predates
2026-08-29 (\hist), so it is orientation, not a controlled comparison, and the comparison it does
support is narrower than it looks: backend, weight format, KV precision and speculative method all
change together between rows.

\begin{table}[t]
\centering\small
\caption{One probe across engines\hist: greedy, 1{,}024-token generation, single stream, median of three
runs, \textbf{decode into a nearly empty KV cache}. vLLM runs tensor-parallel across both cards with
\flag{--enforce-eager}. Power is GPU telemetry plus CPU package plus a fixed platform allowance
(modelled). Context ceiling is the largest window that boots (vLLM) or that prefills and generates
(llama.cpp); the llama.cpp ceilings are first-battery values from \cref{sec:context}, the only
non-historical cells in this table, shown for orientation.}\label{tab:backends}
\shrinktofit{\begin{tabular}{@{}llrrrrrr@{}}
\toprule
backend / weights & speculation & tok/s & speed-up & accept. & avg W & J/tok & ctx ceiling \\
\midrule
vLLM / NVFP4 & none & 12.44 & 1.00$\times$ & --- & 126 & 9.96 & 51{,}200 \\
vLLM / NVFP4 & MTP $n{=}2$ & 28.11 & 2.26$\times$ & 0.709 & 133 & 4.57 & $\approx$49K \\
vLLM / NVFP4 & MTP $n{=}4$ & \textbf{41.21} & \textbf{3.31}$\times$ & 0.559 & 148 & \textbf{3.60} & $\approx$49K \\
vLLM / NVFP4 & DFlash2 (BF16 drafter) & \multicolumn{6}{l}{infeasible: out of memory at every setting tried} \\
\addlinespace
llama.cpp / \Qfour & none & 23.00 & 1.00$\times$ & --- & 207 & 9.00 & 262{,}144 \\
llama.cpp / \Qfour & MTP $n{=}2$ & 32.69 & 1.42$\times$ & 0.728 & 214 & 6.55 & 262{,}144 \\
llama.cpp / \Qfour & DFlash2 $n{=}4$ (4-bit drafter) & 38.51 & 1.67$\times$ & 0.714 & 218 & 5.67 & 262{,}144 \\
\addlinespace
SGLang 0.5.18 & EAGLE & \multicolumn{6}{l}{never started: out of memory initialising the draft worker} \\
\bottomrule
\end{tabular}}
\end{table}

\paragraph{vLLM with NVFP4.} Tensor parallelism splits every layer across both cards, so vLLM balances
memory better than llama.cpp's layer split does. It lost on two other axes. \emph{Context:} at the
highest GPU-memory-utilisation setting that boots (0.97; at 0.99 the server dies within a minute on
16\,GB cards) the nightly build reports a KV pool of 52{,}337 tokens, a working ceiling of 51{,}200,
about one fifth of the model's native window, and MTP takes a little more. An earlier run on the
stable image recorded 98{,}304; the two were never measured back to back, so the ceiling may be
image-dependent. Neither figure was verified with a deep prefill, because that battery was stopped
early. \emph{Fidelity:} NVFP4 scored 85.4/84.1 on HumanEval/HumanEval+ in thinking mode against
90.9/88.4 for \QsixXL{} (\cref{tab:humaneval}); on matched perplexity windows it sits just below the
whole GGUF ladder (\cref{sec:ppl}); and Unsloth's published KL divergence for it on code is 0.026
against roughly 0.002 for their 6-bit dynamic build \citep{unsloth_dynamic3}. None of those gaps is
large, and the HumanEval+ one is inside its interval. Taken together with a five-fold context deficit
they decided the practical question.

\paragraph{The larger speculative multiplier belongs to the slower baseline.} In \cref{tab:backends}
vLLM gains $3.3\times$ from MTP and llama.cpp gains $1.4$--$1.7\times$, yet the best absolute rates
differ by 7\,\%. The \flag{--enforce-eager} setting this configuration needs slows vLLM's plain decode to
12.4\tps, which leaves more for speculation to win back. A speed-up ratio without its baseline says
little. One thing does carry across engines: MTP acceptance at depth 2 is 0.728 on llama.cpp and 0.709
on vLLM, which suggests acceptance is mostly a property of the model's draft head, not of either
implementation.

\paragraph{A separate draft model hits a wall on 16\,GB cards.} SGLang failed at every context length
tried (131K down to 16K) with \texttt{Tried to allocate 1.19\,GiB} while initialising its EAGLE draft
worker. vLLM with the 3.6\,GB BF16 DFlash2 drafter failed on an allocation of the \emph{same size} at
every utilisation setting from 0.78 to 0.97. llama.cpp with a 1.1\,GB 4-bit GGUF drafter worked. So
three stacks were tried, two failed the same way, and the one that succeeded differs in the two ways
one would predict: a quantized drafter, and layer-split placement that leaves per-card headroom where
tensor-parallel replication does not. The matching 1.19\,GiB figure is an observation; no allocator
was instrumented. MTP avoids the problem on both backends because its predictor lives inside the
checkpoint. SGLang produced no measurement of any kind, and nothing here says anything about its
quality: a non-speculative SGLang configuration was not attempted.

\begin{takeaway}
\textbf{Why llama.cpp.} On two unlinked 16\,GB cards it was the only server that reached the native
262K window, the only one that ran a separate drafter, and the one with the wider choice of weight
precisions. vLLM's tensor parallelism balances memory better and its MTP path was the most
energy-efficient configuration measured (3.60\,J/token at depth 0), so this is a statement about this
memory budget and these builds, not about the engines in general.
\end{takeaway}
